\section{Symmetrised three-point function $\propto d^{abc}$: the odd part of the finite terms}\label{sec:dabc}
\noindent\textbf{Result in one paragraph.} The author replaces (A2) by
\begin{align}
\langle\{\rho^{a}(x)\rho^{b}(y)\rho^{c}(z)\}_{\rm sym}\rangle=d^{abc}\,\frac{N_{c}}{(N_{c}^{2}-4)(N_{c}^{2}-1)}\,o(x,y,z)\,,\qquad o(x,y,z)\equiv\langle d^{efg}\{\rho^{e}(x)\rho^{f}(y)\rho^{g}(z)\}_{\rm sym}\rangle\,,\tag{A2$'$}
\end{align}
(normalisation confirmed by the author; $d^{abc}d^{abc}=(N_{c}^{2}-4)(N_{c}^{2}-1)/N_{c}$; $o$ is real and totally symmetric). Then every cubic finite term acquires, in addition to its contact terms (which are unchanged, \texttt{fix\_singlet}), a \emph{symmetrised part} $\propto o$. For every cubic colour structure of the finite terms we give it in colour-singlet form: it is the row's $c$-number factor times $\frac{N_{c}}{(N_{c}^{2}-4)(N_{c}^{2}-1)}\,\Delta\,o$, where $\Delta=\sum_{\rm orderings}C^{abc}d^{abc}$ is a combination of a single new invariant $\mf{G}(a,b,c)=f^{ijk}d^{lmn}U^{il}(a)U^{jm}(b)U^{kn}(c)$ (totally antisymmetric, C-odd) and, in Rows II.I and III.II, of its four-line analogue $\mf{Q}_{d}$. In the fundamental representation
\begin{align}
\mf{G}(a,b,c)={\rm Im}\big[{\rm tr}(V_{a}^{\dagger}V_{b})\,{\rm tr}(V_{b}^{\dagger}V_{c})\,{\rm tr}(V_{c}^{\dagger}V_{a})\big]+{\rm Im}\,{\rm tr}\big(V_{a}^{\dagger}V_{b}V_{c}^{\dagger}V_{a}V_{b}^{\dagger}V_{c}\big)\,,\notag
\end{align}
derived exactly (Fierz) and verified for $N_{c}=2,\dots,6$. \textbf{All thirteen cubic colour structures (Rows I.II $N_{2b}$, I.III, I.V first part, II.I first part, III.II, III.III first part, III.V, III.VI) have a non-zero symmetrised part}; none vanishes identically. With the c.c.\ (Groups I, III), resp.\ with $N_{1a}+N_{1b}$ (Group II), it is real and \emph{odd} under $k\to-k$: it is the odderon-type, $k$-odd part of the single inclusive spectrum. It is C-odd in the projectile ($o$) and in the target ($\mf{G}$, $\mf{Q}_{d}$ change sign under $V\to V^{*}$; for a weak target they start at three gluon exchanges), so it vanishes after the target average for a C-even target ensemble (MV), and it vanishes for $SU(2)$. \textbf{The three-$\rho$ parts of the logarithmic (BFKL, JIMWLK) terms still cancel} for this three-point function: the cancellation is an identity of the classical (totally symmetrised) colour tensors, verified now also for $SU(3)$ and $SU(4)$ and in the $d^{abc}$ projection; this is consistent with the LO term having no three-$\rho$ part (Sec.~\ref{sec:dabc-evol}). Numerical checks: colour algebra $2.4\times10^{-15}$; fundamental forms $2.3\times10^{-15}$; operator level (explicit $SU(3)$ charges, C-odd ensemble with $d$-type three-point function) $7.7\times10^{-16}$; parity exact; evolution cancellation $|{\rm Sym}_{3}|\le2\times10^{-15}$ against termwise $d$-projections up to $44$. Drop-in replacements for the paper and for \texttt{Final\_terms\_fixed} are in Sec.~\ref{sec:dabc-dropin}. Scripts in \texttt{scratchpad/fix/dabc/} (Sec.~\ref{sec:dabc-num}).

\subsection{The assumption and its normalisation}\label{sec:dabc-ass}
\noindent By (A1) (global colour invariance) the symmetrised three-point function $S^{abc}(x,y,z)\equiv\langle\{\rho^{a}(x)\rho^{b}(y)\rho^{c}(z)\}_{\rm sym}\rangle$ ($\{\dots\}_{\rm sym}$: average over the six orderings) is an invariant tensor, $S^{abc}=d^{abc}g_{d}+f^{abc}g_{f}$, with $g_{d}$ real and totally symmetric, $g_{f}$ real and totally antisymmetric in the three points (\texttt{fix\_singlet}, Sec.~``Projectile average''). Contracting with $d^{abc}$ and using $d^{abc}f^{abc}=0$, $d^{abc}d^{abc}=(N_{c}^{2}-4)(N_{c}^{2}-1)/N_{c}$,
\begin{align}
g_{d}(x,y,z)=\frac{N_{c}}{(N_{c}^{2}-4)(N_{c}^{2}-1)}\,o(x,y,z)\,,\qquad g_{f}(x,y,z)=\frac{1}{N_{c}(N_{c}^{2}-1)}\,\langle f^{efg}\{\rho^{e}(x)\rho^{f}(y)\rho^{g}(z)\}_{\rm sym}\rangle\,.\label{eq:dabc-gdgf}
\end{align}
(A2$'$) is therefore equivalent to (A1) together with $g_{f}=0$, and it reduces to the former (A2) when $o=0$ (Gaussian projectile, or any ensemble invariant under the antiunitary map $\rho^{a}\to-\rho^{a}$, \texttt{fix\_singlet}). $o$ is C-odd: under charge conjugation of the projectile, $\rho^{a}t^{a}\to-(\rho^{a}t^{a})^{T}$, the cubic Casimir density $d^{efg}\rho^{e}\rho^{f}\rho^{g}$ changes sign; $o$ is the projectile odderon. A model in which it is non-zero is the weight functional of Jeon and Venugopalan (hep-ph/0503219, Eq.~(17)), $\exp\{-\int_{x}[\rho^{a}\rho^{a}/(2\mu_{A}^{2})-d_{abc}\rho^{a}\rho^{b}\rho^{c}/\kappa_{A}]\}$, whose cubic term ``generates Odderon excitations'' (their abstract); to first order in $1/\kappa_{A}$ it gives (our evaluation by Wick's theorem, their normalisation of $\mu_{A}$, $\kappa_{A}$, classical $\rho$ = Weyl symbol) $S^{abc}=\frac{6\mu_{A}^{6}}{\kappa_{A}}d^{abc}\delta^{(2)}(x-y)\delta^{(2)}(x-z)$, i.e.\ a \emph{local} $o$ (see Sec.~\ref{sec:dabc-local}).\par
\noindent\emph{$SU(2)$:} $d^{abc}\equiv0$, so $o\equiv0$ and the normalisation factor is $0/0$; for $SU(2)$, (A2$'$) is to be read as $S^{abc}=0$ (the only invariant is $f^{abc}$, which (A2$'$) excludes). All results below vanish identically for $SU(2)$ ($\mf{G}=\mf{Q}_{d}=0$).

\subsection{General formula}
\noindent Every cubic finite term has the form $A(k)\,C^{abc}\rho^{a}(p_{1})\rho^{b}(p_{2})\rho^{c}(p_{3})$ (ordered, type ``ord''), or with an anticommutator, $\rho(p_{1})\{\rho(p_{2}),\rho(p_{3})\}$ (A23) or $\{\rho(p_{1}),\rho(p_{2})\}\rho(p_{3})$ (A12); $A(k)$ is the $c$-number factor (row prefactor and kernel), $C$ is real (adjoint $U$, $f^{abc}$). Expanding the anticommutators into orderings, the symmetrised part of the average is $\sum_{\rm orderings}C^{abc}_{\rm seq}S^{abc}(\text{points of seq})$; since $S$ is symmetric under simultaneous permutations of indices and points, and $g_{d}$ totally symmetric,
\begin{align}
\big\langle C^{abc}\rho^{a}(p_{1})\rho^{b}(p_{2})\rho^{c}(p_{3})\big\rangle\big|_{\rm sym}=\frac{N_{c}}{(N_{c}^{2}-4)(N_{c}^{2}-1)}\,\Delta\,o(p_{1},p_{2},p_{3})\,,\qquad\Delta\equiv\sum_{\rm orderings}C^{abc}_{\rm seq}\,d^{abc}\,,\label{eq:dabc-gen}
\end{align}
($\Delta=C^{abc}d^{abc}$ for ``ord'', $\Delta=2C^{abc}d^{abc}$ for A23, A12, with $C$ the coefficient of the first ordering). The full average is
\begin{align}
\langle C\rho\rho\rho\rangle=\frac{N_{c}}{(N_{c}^{2}-4)(N_{c}^{2}-1)}\,\Delta\,o(p_{1},p_{2},p_{3})+[\text{contact terms of \texttt{fix\_singlet}, Eqs.~(s-ord)--(s-a12), unchanged}]\,.\notag
\end{align}
\emph{If $g_{f}\neq0$} (not part of (A2$'$)), the same argument adds $t\,\Phi\,g_{f}(p_{1},p_{2},p_{3})$, $t=1$ (ord), $2$ (A12, A23), with $\Phi=C^{abc}f^{abc}$ the singlet factor of \texttt{fix\_singlet} at \emph{separate} points (verified at the operator level, Sec.~\ref{sec:dabc-num}).

\subsection{The invariants $\mf{G}$ and $\mf{Q}_{d}$}\label{sec:dabc-inv}
\noindent Every cubic colour structure has exactly one $f$; contracting its three $\rho$ indices with $d$ gives a network with one $f$, one $d$ and chains of Wilson lines. Chains joining two legs of $f$ to two legs of $d$ directly give $f^{abe}d^{abe'}=0$ (e.g.\ $f^{hbc}U^{ah}(w)\{\rho^{b},\rho^{c}\}$, $U^{ac'}(w')U^{ac}(w)f^{cbd}\rho^{c'}\rho^{b}\rho^{d}$); the others reduce, with $U^{T}U=1$, to (notation of \texttt{fix\_singlet}, $(T^{e})_{bc}=-if^{ebc}$, and $(D^{e})_{bc}\equiv d^{ebc}$)
\begin{align}
\mf{G}(a,b,c)&\equiv f^{ijk}\,d^{lmn}\,U^{il}(a)\,U^{jm}(b)\,U^{kn}(c)\,,\notag\\
\mf{Q}_{d}(a,b;c,e)&\equiv i\,{\rm Tr}\{T^{h}\,U^{\dagger}(a)U(b)\,D^{h}\,U^{\dagger}(c)U(e)\}=f^{hij}d^{hkl}\,[U^{\dagger}(a)U(b)]^{jk}\,[U^{\dagger}(c)U(e)]^{li}\,.\label{eq:dabc-GQ}
\end{align}
\textbf{Properties} (all verified numerically, Sec.~\ref{sec:dabc-num}): $\mf{G}$ is real and totally antisymmetric; $\mf{G}(a,a,c)=0$; $\mf{Q}_{d}$ is real, $\mf{Q}_{d}(a,b;c,e)=-\mf{Q}_{d}(e,c;b,a)$, $\mf{Q}_{d}(a,a;c,e)=\mf{Q}_{d}(a,b;c,c)=0$, $\mf{Q}_{d}(a,b;c,a)=-\mf{G}(a,b,c)$, and $\mf{Q}_{d}(a,b;b,c)=\mf{G}(a,b,c)|_{U\to U^{T}}$ (i.e.\ $V\to V^{\dagger}$; not reducible to $\mf{G}$ itself). Both are invariant under $U(p)\to RU(p)R'$ for all $p$ (global rotations $R,R'$), hence vanish for constant $U$, and both are \emph{C-odd}: under $V\to V^{*}$, $U\to\eta U\eta$ with $\eta={\rm diag}(\pm1)$, and $\eta_{a}\eta_{b}\eta_{c}f^{abc}=-f^{abc}$, $\eta_{a}\eta_{b}\eta_{c}d^{abc}=+d^{abc}$, so $\mf{G}\to-\mf{G}$, $\mf{Q}_{d}\to-\mf{Q}_{d}$ (whereas $\mf{D},\mf{F},\mf{Q}$ are C-even). For a weak target, $V_{p}=e^{i\epsilon\alpha_{p}^{a}t^{a}}$, $\mf{G}$ and every $\Delta$ below are $O(\epsilon^{3})$: the odd part needs three gluon exchanges with the target (a C-odd, $d^{abc}$ exchange). For $SU(2)$, $\mf{G}=\mf{Q}_{d}=0$.

\noindent\textbf{Fundamental-trace forms (exact).} With $U^{ab}(p)=2{\rm tr}(t^{a}V_{p}t^{b}V_{p}^{\dagger})$ and $A\equiv V_{a}^{\dagger}V_{b}$, $B\equiv V_{c}^{\dagger}V_{e}$,
\begin{align}
\mf{Q}_{d}(a,b;c,e)&=-{\rm Im}\big[{\rm tr}A\,{\rm tr}B\,{\rm tr}(A^{\dagger}B^{\dagger})\big]-{\rm Im}\,{\rm tr}\big(ABA^{\dagger}B^{\dagger}\big)\,,\label{eq:dabc-Qfund}\\
\mf{G}(a,b,c)=-\mf{Q}_{d}(a,b;c,a)&={\rm Im}\big[{\rm tr}(V_{a}^{\dagger}V_{b})\,{\rm tr}(V_{b}^{\dagger}V_{c})\,{\rm tr}(V_{c}^{\dagger}V_{a})\big]+{\rm Im}\,{\rm tr}\big(V_{a}^{\dagger}V_{b}V_{c}^{\dagger}V_{a}V_{b}^{\dagger}V_{c}\big)\,.\label{eq:dabc-Gfund}
\end{align}
\emph{Derivation.} $\sum_{h}f^{hij}d^{hkl}=-2i\,{\rm tr}([t^{i},t^{j}]\{t^{k},t^{l}\})$ (completeness; the $1/N_{c}$ term drops since ${\rm tr}[t^{i},t^{j}]=0$), and $\sum_{j}t^{j}[U^{\dagger}(a)U(b)]^{jk}=At^{k}A^{\dagger}$, $\sum_{i}t^{i}[U^{\dagger}(c)U(e)]^{li}=B^{\dagger}t^{l}B$, so $\mf{Q}_{d}=-2i\sum_{k,l}{\rm tr}\big([B^{\dagger}t^{l}B,At^{k}A^{\dagger}]\{t^{k},t^{l}\}\big)$. The sums over $k$ and $l$ with $\sum_{k}{\rm tr}(Pt^{k}Qt^{k})=\frac12{\rm tr}P\,{\rm tr}Q-\frac{1}{2N_{c}}{\rm tr}(PQ)$ and $\sum_{k}{\rm tr}(Pt^{k}){\rm tr}(Qt^{k})=\frac12{\rm tr}(PQ)-\frac{1}{2N_{c}}{\rm tr}P\,{\rm tr}Q$ give, after all $1/N_{c}$ terms cancel, $\sum_{k,l}(\dots)=-\frac{i}{2}\big[{\rm Im}\,Z_{1}+{\rm Im}\,Z_{2}\big]$ with $Z_{1}={\rm tr}A\,{\rm tr}B\,{\rm tr}(A^{\dagger}B^{\dagger})$, $Z_{2}={\rm tr}(ABA^{\dagger}B^{\dagger})$; hence Eq.~\eqref{eq:dabc-Qfund}. The coefficients are $N_{c}$-independent; the same result came out of an unbiased fit to all trace monomials for $N_{c}=4,\dots,7$ and was checked for $N_{c}=2,\dots,6$ (Sec.~\ref{sec:dabc-num}). These are the standard C-odd (odderon) combinations: ${\rm Im}$ of the product of three dipoles around the triangle $a,b,c$ and ${\rm Im}$ of the sextupole; both are odd under the exchange of two points.

\subsection{The symmetrised part of every cubic structure}\label{sec:dabc-table}
\noindent $\Delta$ for each cubic colour structure (labels and $\rho$ orderings of \texttt{fix\_singlet}; the structures are those of \texttt{Final\_terms\_fixed}, identical to \texttt{Final\_terms\_final}, \texttt{Final\_terms\_checked} and the paper appendix), with arguments ordered by the list $x,y,x',y',z,w,w'$. ``Local'' is $\Delta$ with all three $\rho$ points set equal (relevant for a local $o$, Sec.~\ref{sec:dabc-local}).
\begin{center}
\renewcommand{\arraystretch}{1.35}
\small\begin{tabular}{|l|l|l|l|}\hline
structure & $\rho$ points, type & $\Delta=\sum C^{abc}d^{abc}$ & local \\\hline
I.II, $N_{2b}$ & $x',x,y$; A23 & $2\big[\mf{G}(x,y,w')-\mf{G}(x,y,x')\big]$ & $0$\\
I.III, $C_{D}=f^{abd}\mathcal{C}_{\rm III}(z)$ & $x',x,y$; ord & $\mf{G}(x,y,x')-\mf{G}(x,y,w')+\mf{G}(x,x',z)+\mf{G}(x,z,w')$ & $\mf{G}(x,z,w')$\\
I.III, $C_{A}=f^{abd}\mathcal{C}_{\rm III}(w)$ & $x',x,y$; ord & $\mf{G}(x,y,x')-\mf{G}(x,y,w')+\mf{G}(x,x',w)+\mf{G}(x,w,w')$ & $\mf{G}(x,w,w')$\\
I.V, $N_{1}$ & $x',y,y'$; A23 & $-2\big[\mf{G}(x,x',z)+\mf{G}(x,z,w')+\mf{G}(y,x',y')+\mf{G}(y,y',w')\big]$ & $-2\mf{G}(x,y,z)-2\mf{G}(x,z,w')$\\
II.I, $N_{1a}$ & $x',y,x$; A12 & $2\big[\mf{G}(x,z,w')-\mf{G}(x,y,x')-\mf{Q}_{d}(z,y;x',w)\big]$ & $2\mf{G}(x,z,w')-2\mf{Q}_{d}(z,x;x,w)$\\
II.I, $N_{1b}$ & $x,x',y$; A23 & $2\big[\mf{G}(x,z,w)-\mf{G}(x,y,x')-\mf{Q}_{d}(z,y;x',w')\big]$ ($=N_{1a}|_{w\leftrightarrow w'}$) & $2\mf{G}(x,z,w)-2\mf{Q}_{d}(z,x;x,w')$\\
III.II & $y,x',x$; ord & $\mf{G}(x,y,x')+\mf{G}(x,x',z)+\mf{Q}_{d}(z,y;x',w)$ & $\mf{Q}_{d}(z,x;x,w)$\\
III.III$_{1}$, $\mathcal{L}(y')$ & $x',y,x$; A23 & $2\big[\mf{G}(x,y,x')-\mf{G}(x,y,y')-\mf{G}(x',z,w)+\mf{G}(y',z,w)\big]$ & $2\mf{G}(y',z,w)-2\mf{G}(x,z,w)$\\
III.III$_{1}$, $\mathcal{L}(z)$ & $x',y,x$; A23 & $2\big[\mf{G}(x,y,x')-\mf{G}(x,y,z)-\mf{G}(x',z,w)\big]$ & $-2\mf{G}(x,z,w)$\\
III.V, $\mathcal{L}_{\rm V}(y')$ & $x',x,y$; ord & $-\mf{G}(x,y,x')+\mf{G}(x,y,y')-\mf{G}(x,x',w)+\mf{G}(x,y',w)$ & $\mf{G}(x,y',w)$\\
III.V, $\mathcal{L}_{\rm V}(z)$ & $x',x,y$; ord & $-\mf{G}(x,y,x')+\mf{G}(x,y,z)-\mf{G}(x,x',w)+\mf{G}(x,z,w)$ & $\mf{G}(x,z,w)$\\
III.VI, $\mathcal{L}(y')$ & $x',y,x$; ord & $-\mf{G}(x,y,x')+\mf{G}(x,y,y')+\mf{G}(y,x',z)-\mf{G}(y,y',z)$ & $-\mf{G}(x,y',z)$\\
III.VI, $\mathcal{L}(z)$ & $x',y,x$; ord & $-\mf{G}(x,y,x')+\mf{G}(x,y,z)+\mf{G}(y,x',z)$ & $0$\\\hline
\end{tabular}
\end{center}
The $\mathcal{L}(z)$ forms are the $\mathcal{L}(y')$ forms with $y'\to z$ (using $\mf{G}(a,a,c)=0$), and $C_{A}$ is $C_{D}$ with $z\to w$. All thirteen $\Delta$ are non-zero for generic points (typical values $0.2$--$6$ for random $SU(3)$, $SU(4)$, $SU(5)$ lines, Sec.~\ref{sec:dabc-num}). The quadratic rows (I.II $N_{1}$, I.V second part, II.I second part, II.IV, III.III second part) and the DGLAP and running-coupling terms are unchanged. \emph{Pairs:} the colour identities~(colid) of the soft-remainder pairs are operator identities, so the symmetrised parts pair exactly as the contact terms do; e.g.\ $\Delta_{\rm III.V}(\mathcal{L}_{\rm V}(y'))|_{y'\to w'}=-\Delta_{\rm I.III}(C_{A})$ and $\Delta_{\rm III.VI}(\mathcal{L}(y'))|_{y'\to w',\,x\leftrightarrow y}=+\Delta_{\rm I.III}(C_{D})$ (checked numerically, $9\times10^{-16}$).

\noindent\textbf{Row factors.} From the row prefactors of \texttt{Final\_terms\_fixed} (equivalently: the contact line coefficient of \texttt{fix\_singlet} divided by $c\,i$), in units of $\mathcal{N}_{\rm f}=\frac{1}{(2\pi)^{3}}\frac{g^{4}}{16\pi^{5}}\frac{1}{k^{+}}$ ($=\mathcal{C}_{\rm f}$ of the paper), the $c$-number factor $A$ of each cubic structure is: I.II $N_{2b}$: $-\frac{1}{4\pi}\mathfrak{K}^{\rm II}_{2b}$; I.III: $2i\,\mathfrak{K}^{\rm III}_{z}$ ($C_{D}$), $2i\,\mathfrak{K}^{\rm III}_{w}$ ($C_{A}$); I.V $N_{1}$: $\frac{i}{2}\mathfrak{K}^{\rm V}_{1}$; II.I: $-i\,\mathfrak{K}_{\mathcal{T}}(w,w')$ ($N_{1a}$), $+i\,\mathfrak{K}_{\mathcal{T}}(w',w)$ ($N_{1b}$); III.II: $-2i\,\mathfrak{K}_{\mathcal{T}}(w,w')$; III.III first part: $-i[-\mathfrak{K}_{\widehat{\mathcal{E}}}+\mathfrak{K}_{\Psi}+\mathfrak{K}^{(y')}_{\delta}]$ ($\mathcal{L}(y')$), $-i\,\mathfrak{K}^{(z)}_{\delta}$ ($\mathcal{L}(z)$); III.V and III.VI: $-2i[-\mathfrak{K}_{\widehat{\mathcal{E}}}+\mathfrak{K}_{\Psi}+\mathfrak{K}^{(y')}_{\delta}]$ ($\mathcal{L}(y')$), $-2i\,\mathfrak{K}^{(z)}_{\delta}$ ($\mathcal{L}(z)$); times the phase of the row. (Spot checks against the printed prefactors: I.III $\frac{ig^{4}}{8\pi^{5}}$, II.I $\mp\frac{ig^{4}}{16\pi^{5}}$, III.II $-\frac{ig^{4}}{8\pi^{5}}$, III.III $\frac{ig^{4}}{16\pi^{5}}$, III.V $\frac{ig^{4}}{8\pi^{5}}$.)

\noindent\textbf{Singlet form of the odd part.} With the kernels of \texttt{fix\_singlet} / the paper,
\begin{align}
\frac{d^{3}N_{\rm odd}}{d^{3}k}\bigg|_{\rm singlet}&=\frac{\mathcal{N}_{\rm f}}{N_{c}^{2}-1}\,\frac{N_{c}}{N_{c}^{2}-4}\Bigg\{\int_{w,w'}e^{-ik\cdot(w-w')}\int_{x,y,x'}\frac{1}{2\pi}\,\mathfrak{K}^{\rm II}_{2b}\Big[\mf{G}(x,y,x')-\mf{G}(x,y,w')\Big]\,o(x,y,x')\notag\\
&\quad+\int_{w,w'}e^{-ik\cdot(w-w')}\int_{x,y,z,x'}2i\Big\{\mathfrak{K}^{\rm III}_{z}\Big[\mf{G}(x,y,x')-\mf{G}(x,y,w')+\mf{G}(x,x',z)+\mf{G}(x,z,w')\Big]\notag\\
&\qquad\qquad+\mathfrak{K}^{\rm III}_{w}\Big[\mf{G}(x,y,x')-\mf{G}(x,y,w')+\mf{G}(x,x',w)+\mf{G}(x,w,w')\Big]\Big\}\,o(x,y,x')\notag\\
&\quad-\int_{x,z,y,y',x',w'}e^{-ik\cdot(z-w')}\,i\,\mathfrak{K}^{\rm V}_{1}\Big[\mf{G}(x,x',z)+\mf{G}(x,z,w')+\mf{G}(y,x',y')+\mf{G}(y,y',w')\Big]\,o(y,x',y')\notag\\
&\quad-\int_{x,y,x',z,w,w'}e^{-ik\cdot(w-w')}\,2i\,\mathfrak{K}_{\mathcal{T}}(w,w')\Big[\mf{G}(x,y,x')+\mf{G}(x,x',z)+\mf{Q}_{d}(z,y;x',w)\Big]\,o(x,y,x')\notag\\
&\quad-\int_{w,x',y',z,y,x}e^{-ik\cdot(w-y')}\,2i\,\Big\{\Big[-\mathfrak{K}_{\widehat{\mathcal{E}}}+\mathfrak{K}_{\Psi}+\mathfrak{K}^{(y')}_{\delta}\Big]\Big[\mf{G}(x,y,x')-\mf{G}(x,y,y')-\mf{G}(x',z,w)+\mf{G}(y',z,w)\Big]\notag\\
&\qquad\qquad+\mathfrak{K}^{(z)}_{\delta}\Big[\mf{G}(x,y,x')-\mf{G}(x,y,z)-\mf{G}(x',z,w)\Big]\Big\}\,o(x,y,x')\notag\\
&\quad-\int_{w,y',z,y,x,x'}e^{-ik\cdot(w-y')}\,2i\,\Big\{\Big[-\mathfrak{K}_{\widehat{\mathcal{E}}}+\mathfrak{K}_{\Psi}+\mathfrak{K}^{(y')}_{\delta}\Big]\Big[-\mf{G}(x,y,x')+\mf{G}(x,y,y')-\mf{G}(x,x',w)+\mf{G}(x,y',w)\Big]\notag\\
&\qquad\qquad+\mathfrak{K}^{(z)}_{\delta}\Big[-\mf{G}(x,y,x')+\mf{G}(x,y,z)-\mf{G}(x,x',w)+\mf{G}(x,z,w)\Big]\Big\}\,o(x,y,x')\notag\\
&\quad-\int_{w,y',z,y,x,x'}e^{-ik\cdot(w-y')}\,2i\,\Big\{\Big[-\mathfrak{K}_{\widehat{\mathcal{E}}}+\mathfrak{K}_{\Psi}+\mathfrak{K}^{(y')}_{\delta}\Big]\Big[-\mf{G}(x,y,x')+\mf{G}(x,y,y')+\mf{G}(y,x',z)-\mf{G}(y,y',z)\Big]\notag\\
&\qquad\qquad+\mathfrak{K}^{(z)}_{\delta}\Big[-\mf{G}(x,y,x')+\mf{G}(x,y,z)+\mf{G}(y,x',z)\Big]\Big\}\,o(x,y,x')\Bigg\}+{\rm c.c.}\notag\\
&\quad-\frac{\mathcal{N}_{\rm f}}{N_{c}^{2}-1}\,\frac{N_{c}}{N_{c}^{2}-4}\int_{x,y,z,w,x',w'}4\sin\big(k\cdot(w-w')\big)\,\mathfrak{K}_{\mathcal{T}}(w,w')\Big[\mf{G}(x,z,w')-\mf{G}(x,y,x')-\mf{Q}_{d}(z,y;x',w)\Big]\,o(x,y,x')\,.\label{eq:dabc-odd}
\end{align}
Lines: I.II $N_{2b}$; I.III; I.V first part; III.II; III.III first part; III.V; III.VI ($+$c.c.\ for all of these); last line: II.I first part, $N_{1a}+N_{1b}$ (no c.c.), obtained from $-i\,e^{-ik\cdot(w-w')}\mathfrak{K}_{\mathcal{T}}(w,w')\Delta_{1a}+i\,e^{-ik\cdot(w-w')}\mathfrak{K}_{\mathcal{T}}(w',w)\Delta_{1b}$ by $w\leftrightarrow w'$ in the second term ($\Delta_{1b}=\Delta_{1a}|_{w\leftrightarrow w'}$). In the III.III--III.VI lines the pairs $\mathfrak{K}^{(y')}_{\delta}$, $\mathfrak{K}^{(z)}_{\delta}$ stay under one $y'$ integral, as for the contact terms. The \emph{full} singlet form of a cubic row is its \texttt{fix\_singlet} (contact) form plus its line of Eq.~\eqref{eq:dabc-odd}.

\subsection{Parity, reality, target average}\label{sec:dabc-parity}
\noindent $\Delta$ and $o$ are real. For Groups I and III the term plus its c.c.\ has the symmetrised part $2\,{\rm Re}[A(k)]\,\frac{N_{c}}{(N_{c}^{2}-4)(N_{c}^{2}-1)}\Delta\,o$. All kernels depend on $k$ through $e^{-ik\cdot r}$, $\mathcal{E}$, $\mathcal{J}$, $\widehat{\mathcal{E}}$, $\Psi$, $\pi\delta(s)-i{\rm P}\frac1s$, all of the form $\int dt\,(\text{real})\,e^{-ist}$, so $f(-s)=f(s)^{*}$; the cubic rows have an explicit $i$ (in I.II $N_{2b}$ the kernel $\mathfrak{K}^{\rm II}_{2b}$ is imaginary, the Fourier transform of the odd real $\mathcal{R}^{i}$). Hence $A(-k)=-A(k)^{*}$, ${\rm Re}A(-k)=-{\rm Re}A(k)$: \textbf{the symmetrised part is real and odd in $k$, pointwise in the coordinates}. For II.I the integrand of the last line of~\eqref{eq:dabc-odd} is manifestly real and odd. (Checked numerically with the closed forms of $\mathcal{E}$, $\widehat{\mathcal{E}}$, $\Psi$, $\mathcal{J}$, for I.III $C_{A}$, III.V $\mathcal{L}_{\rm V}(y')$ and II.I $N_{1a}+N_{1b}$, Sec.~\ref{sec:dabc-num}.) The contact terms remain even; hence, under (A1)+(A2$'$), \emph{the $k$-odd part of the finite terms is exactly Eq.~\eqref{eq:dabc-odd}}, and the $k$-even part is that of \texttt{fix\_singlet}.\par
\noindent\emph{When does it survive?} (i) It is C-odd in the projectile ($o$) and in the target ($\mf{G},\mf{Q}_{d}$). After an average over a C-even target ensemble (e.g.\ Gaussian MV target) it vanishes; it survives configuration by configuration (event-by-event odd harmonics) or for a target with a C-odd (odderon) component. (ii) All kernels are invariant under the simultaneous rotation by $\pi$ of all coordinates and of $k$ ($r\to-r$, $k\to-k$). Therefore, if the averaged product $\langle\Delta\rangle_{\rm target}\,o$ is invariant under the reflection of all points through a common centre, the odd part integrates to zero; a non-zero result needs a configuration (or an impact-parameter geometry) that breaks this symmetry. (iii) For a weak target it is $O(\epsilon^{3})$ in the target field (three-gluon, $d$-type exchange), and $O(o)$ in the projectile: projectile odderon $\times$ target odderon.

\subsection{Local three-point function}\label{sec:dabc-local}
\noindent For a local $o$ (Jeon--Venugopalan type, $o\propto\delta^{(2)}(x-y)\delta^{(2)}(x-z)$) only the ``local'' column of the table enters: I.II $N_{2b}$ and III.VI $\mathcal{L}(z)$ drop out, the other structures do not. \textbf{Not checked:} whether the kernels are integrable when all three $\rho$ points coincide (\texttt{fix\_singlet} checked only pairwise merging, for two kernels); with a local $o$ this has to be power-counted row by row.

\subsection{Three-$\rho$ parts of the BFKL and JIMWLK terms}\label{sec:dabc-evol}
\noindent \textbf{Statement in the files.} \texttt{Evolution\_three\_rho\_checked} (Sec.~``Cancellation of all three $\rho$'s'') and the paper (Sec.~``Terms cubic in $\rho$'') state that for each logarithm the \emph{totally symmetrised} three-$\rho$ part of the sum of all logarithmic three-$\rho$ terms (genuine ones and the three-$\rho$ remainders of the reordered four-$\rho$ terms) vanishes identically, for each of the twelve classes (kernel $\alpha,\alpha',\beta,\beta'$ $\times$ two, three, four Wilson lines). The check recorded in \texttt{Evolution\_three\_rho\_checked} was done with $SU(2)$ charges (the $SU(3)$ log of that check in \texttt{scratchpad/ev\_3rho} is empty), for which $d^{abc}=0$, so a $d$-type piece could not have been seen there; the paper states $SU(2)$ and $SU(3)$.\par
\noindent\textbf{Re-check with $SU(3)$, $SU(4)$ and the $d$-projection} [\texttt{ev\_generic.py}]. At distinct points the charges commute, so the symmetrised part of a bracket is the classical (commuting) colour tensor. For each of the twelve cancellation brackets as written, and for the collected list of all rows regrouped by kernel class, the colour tensors were evaluated with random $SU(3)$ and $SU(4)$ Wilson lines at generic points: the summed, totally symmetrised tensor (all colour components) is $\le2\times10^{-15}$, and its $d^{abc}$ projection $\le6\times10^{-15}$, while the sums of the moduli of the termwise $d$-projections are $1.6$--$32$ per bracket ($10$--$44$ per kernel class); the two-Wilson-line classes have vanishing $d$-projection term by term ($f^{abc}$ with two indices on $\rho$'s). Since the vanishing holds for the whole tensor, it holds for \emph{any} symmetrised three-point function: $d^{abc}o$ under (A2$'$), and also an $f^{abc}g_{f}$ part. Contact configurations need no separate treatment: the identity holds for arbitrary Wilson lines at the points, hence also when points coincide, and $\langle\{\rho\rho\rho\}_{\rm sym}\rangle$ is given by (A2$'$) at coincident points as well. \textbf{Result: with $\langle\rho\rho\rho\rangle_{\rm sym}\neq0$ the symmetrised three-$\rho$ parts of the logarithmic terms still cancel; there is no $\ln$-enhanced term $\propto o$}. The two-$\rho$ remainders (part of the BFKL and JIMWLK terms) and the one-$\rho$ remainders (zero average) are not affected.\par
\noindent\textbf{Why this is consistent.} The LO spectrum is quadratic in $\rho$; under (A1) it depends on the projectile only through $n$, which is C-even. The logarithmic terms at $O(g^{4})$ are its one-step evolution: in $\ln(\vee/k^{+})$ the evolution of $n$ (dilute projectile: linear BFKL, Eq.~(BFKLrhorho)), in $\ln(k^{+}/\Lambda)$ the JIMWLK evolution of the target operator at fixed $n$. A $\ln$-enhanced term $\propto o$ would mean that $\partial_{Y}n$ contains $o$ (or that the target evolution of the LO operator produces a three-$\rho$ projectile structure). Both are excluded: the evolution Hamiltonians are C-invariant, so the C-even $n$ cannot mix with the C-odd $o$; and the LO operator has no three-$\rho$ term whose evolution could appear. In addition, an $o$-term would have to come with a C-odd target structure ($\mf{G}$, $\mf{Q}_{d}$), whereas the JIMWLK term is $H_{\rm JIMWLK}$ acting on the C-even LO operator. The $o$-dependence of the spectrum starts with the non-logarithmic $g^{4}\rho^{3}$ terms of Eq.~\eqref{eq:dabc-odd} (the ``first saturation correction''); its rapidity evolution (evolution of $o$, the projectile odderon, and of the C-odd target structures) is of higher order, $\alpha_{s}\ln\times g^{4}\rho^{3}$, and is not part of this calculation. The cancellation found is therefore exactly what is required.

\subsection{Literature}\label{sec:dabc-lit}
\noindent Read in this session from the arXiv sources (export.arxiv.org); only what was read is quoted.
\begin{itemize}
\item McLerran, Skokov, arXiv:1611.09870 (NPA 959 (2017) 83): in the classical Yang--Mills dilute--dense calculation the leading order spectrum is symmetric configuration by configuration, $n^{(2)}(k)=n^{(2)}(-k)$; the first odd contribution is $n^{(3)}$, cubic in the projectile charge $\rho_{1}$, with the odd part $\propto{\rm Im}\{\dots f^{abc}\Omega^{a}\Omega^{b}\Omega^{c*}\}$ (their equation labelled \texttt{Eq:FinalAssymitry}); conclusions: the odd harmonics ``in single particle inclusive process \dots average out to zero for a Gaussian ensemble because they are proportional to $\rho_{1}^{3}$. Essentially they are defined by odderon exchanges.'' \emph{Agreement:} LO even, odd part cubic in $\rho$ and odd configuration by configuration, zero for a Gaussian projectile ($o=0$). Our result adds: with non-commuting charges this holds for the symmetrised part only (contact terms are even, \texttt{fix\_singlet}), and with (A2$'$) the odd part survives the projectile average, $\propto o$, in the explicit form~\eqref{eq:dabc-odd}. A term-by-term comparison with that equation (different formalism: classical fields, momentum space) was not made.
\item Kovchegov, Skokov, arXiv:1802.08166 (PRD 97 (2018) 094021): for single-gluon production ``it is impossible to have a cross section contribution which is odd under $k\to-k$ in the case of unpolarized initial and final states: \dots\ if we take the interference of a single $M_{1}$ with a single $M_{3}$ contribution \dots\ the color averaging would give zero'' (their Sec.\ on the odd-harmonic part at order $\alpha_{s}^{4}$, MV-type averaging); odd harmonics of the two-gluon correlation ``require three scatterings in the projectile and three scatterings in the target''. \emph{Consistent}: our odd part vanishes after the average over a C-even (MV) target or a $\pi$-rotation-symmetric ensemble, and over a projectile with $o=0$; it is cubic in the projectile and $O(\epsilon^{3})$ in a weak target field.
\item Kovner, Lublinsky, Skokov, arXiv:1612.07790: in double inclusive production at leading power the antisymmetric part is ``saturated by a charge conjugation odd `condensate', i.e.\ the Odderon'' of the target, odd under $S\to S^{\dagger}$, which ``averages to zero in the MV model''. \emph{Consistent} with the C-oddness of $\mf{G}$, $\mf{Q}_{d}$ in the target.
\item Jeon, Venugopalan, hep-ph/0503219 (PRD 71 (2005) 125003): weight functional with the cubic term $d_{abc}\rho^{a}\rho^{b}\rho^{c}/\kappa_{A}$ (their Eq.~(17)), which generates Odderon excitations; a concrete (local) model for $o$, Sec.~\ref{sec:dabc-ass}.
\item Li, Skokov, arXiv:2104.01879 (first saturation correction to single inclusive production; abstract only read): not compared.
\end{itemize}

\subsection{Numerical checks}\label{sec:dabc-num}
\noindent Scripts in \texttt{scratchpad/fix/dabc/} (\texttt{sun.py}: $SU(N)$ generators, $f$, $d$, random $V$ and $U^{ab}=2{\rm tr}(t^{a}Vt^{b}V^{\dagger})$; \texttt{invariants.py}; \texttt{dred.py}: chain reduction of the $d$-contracted structures; \texttt{dstructs.py}: the table; logs \texttt{*.log}). The colour structures are parsed from the TeX source by the parser of \texttt{fix\_singlet} (\texttt{scratchpad/fix/singlet/rows.py}); they occur verbatim in \texttt{Final\_terms\_fixed} and \texttt{Final\_terms\_checked} (\texttt{check\_sources.py}).
\begin{itemize}
\item \emph{Colour algebra} [\texttt{check\_delta.py}]: $\Delta$ of the table vs direct contraction $\sum_{\rm seq}C^{abc}_{\rm seq}d^{abc}$ of the parsed structures, $SU(3)$, $SU(4)$, $SU(5)$, three sets of random Wilson lines: largest relative deviation $2.4\times10^{-15}$; canonical (ordered-argument) forms $3\times10^{-15}$ [\texttt{canon.py}].
\item \emph{Invariants} [\texttt{t\_inv.py}, \texttt{check\_fund.py}, \texttt{fund.py}]: antisymmetry, coincident-point rules, $\mf{Q}_{d}(a,b;c,a)=-\mf{G}(a,b,c)$, $\mf{Q}_{d}(a,b;b,c)=\mf{G}|_{U\to U^{T}}$, global invariance ($\le2.4\times10^{-15}$); C-oddness under $V\to V^{*}$ (exact); fundamental forms~\eqref{eq:dabc-Qfund}, \eqref{eq:dabc-Gfund}: $\le2.3\times10^{-15}$ for $N_{c}=2,\dots,6$, 20 samples each; the unbiased least-squares fit to all 24 trace monomials gives the same coefficients for $N_{c}=4,5,6,7$ (fit residual $\le6\times10^{-15}$).
\item \emph{Operator level} [\texttt{check\_operator\_d.py}]: charges = generators on three lattice sites (explicit matrices), globally colour-invariant density matrix $e^{-H}/{\rm Tr}e^{-H}$ (ensembles of \texttt{fix\_singlet}); for every cubic structure and every assignment of its $\rho$ points to sites (random complex weights; coincident sites = contacts): exact ordered $\langle C\rho\rho\rho\rangle$ vs [contact terms of \texttt{fix\_singlet}] $+\frac{N_{c}}{(N_{c}^{2}-4)(N_{c}^{2}-1)}\Delta\,o$ with $o$ computed from the state. $SU(3)$ fundamental charges (C-odd, $\max|o|=1.11$, $f$-part $4\times10^{-18}$; $S^{abc}$ equals the (A2$'$) form to $3.5\times10^{-17}$): largest relative deviation $7.7\times10^{-16}$, the $d$-parts ranging from about $1\%$ to more than $100\%$ of the contact parts. Same with a chiral term $0.6f^{abc}\rho_{1}^{a}\rho_{2}^{b}\rho_{3}^{c}$ in $H$ ($g_{f}\neq0$, (A2$'$) violated): deviation up to $6\times10^{-2}$, removed ($\le2.7\times10^{-15}$) by adding $t\,\Phi\,g_{f}$. C-even ensemble: $o=0$, deviation $1.9\times10^{-15}$ (reduces to \texttt{fix\_singlet}).
\item \emph{Parity} [\texttt{check\_parity.py}]: $f(-s)=f(s)^{*}$ for $\mathcal{E},\widehat{\mathcal{E}},\Psi,\mathcal{J}$ (closed forms; $\mathcal{E}$ also by quadrature, $6\times10^{-17}$); for I.III $C_{A}$, III.V $\mathcal{L}_{\rm V}(y')$ (+c.c.) and II.I $N_{1a}+N_{1b}$ (integrand symmetrised under $w\leftrightarrow w'$), at 50 random configurations: symmetrised part real and odd under $k\to-k$ to rounding.
\item \emph{Weak target} [\texttt{check\_dilute.py}]: $\mf{G}$ and all thirteen $\Delta$ scale as $\epsilon^{3}$ (ratios $v(\epsilon)/v(\epsilon/2)\to8$; for II.I $N_{1a}$ and III.II with visible $O(\epsilon^{4})$ corrections).
\item \emph{Evolution terms} [\texttt{ev\_generic.py}]: Sec.~\ref{sec:dabc-evol}.
\end{itemize}

\subsection{Questions for the author}\label{sec:dabc-q}
\begin{itemize}
\item[(Q1)] (A2$'$) sets the $f^{abc}$ part $g_{f}$ of the symmetrised three-point function to zero. By (A1) alone $S^{abc}=d^{abc}g_{d}+f^{abc}g_{f}$, and C-evenness does not remove $g_{f}$ (\texttt{fix\_singlet}). Is $g_{f}=0$ intended (e.g.\ by an additional symmetry of the projectile)? If not, the odd part acquires $t\,\Phi\,g_{f}$ per structure, with $\Phi$ of \texttt{fix\_singlet} at separate points (verified).
\item[(Q2)] For $SU(2)$ the normalisation of (A2$'$) is $0/0$; we read (A2$'$) as $S^{abc}=0$ there. Agreed?
\item[(Q3)] Which model for $o$ is intended? For a local $o$ (Jeon--Venugopalan) two structures drop out and the integrability of the kernels at the triple coincidence of the $\rho$ points has to be checked; this was not done.
\item[(Q4)] The odd part is C-odd in the target: it vanishes after a C-even target average and survives event by event or with a target odderon. Is the result meant event by event (odd harmonics in correlations), or with a target odderon? This decides how the paper should phrase ``the spectrum is odd''.
\item[(Q5)] Presentation: should Eq.~\eqref{eq:dabc-odd} be part of the main result (Sec.~``Average over the projectile'') or an appendix? The drop-ins below put it in the main text, after the contact-term forms.
\end{itemize}

\subsection{Drop-in replacements}\label{sec:dabc-dropin}
\begingroup
\def\v#1{\mathbf{#1}}\providecommand{\Kv}{\mathcal{K}}\renewcommand{\eqref}[1]{(\texttt{\detokenize{#1}})}\renewcommand{\ref}[1]{\texttt{\detokenize{#1}}}\renewcommand{\cite}[1]{[\texttt{\detokenize{#1}}]}\providecommand{\pending}[1]{\textcolor{red}{[#1]}}
\noindent Labels, references and citations are printed as their keys here; in the source they are the paper's. $\mathcal{C}_{\rm f}=\mathcal{N}_{\rm f}$.

\subsubsection*{(P1) Paper, \texttt{sec\_finite.tex}, Sec.~\ref{sec:fin_singlet}: item (A2) and the paragraph after the itemize}
\noindent Replace the item (A2) by:\par\smallskip
\noindent\fbox{\parbox{0.97\textwidth}{\small\begin{itemize}
\item[(A2)] \emph{Symmetrised three-point function of $d$-type:} with $S^{abc}(\v x,\v y,\v z)\equiv\langle\{\rho^a(\v x)\rho^b(\v y)\rho^c(\v z)\}_{\rm sym}\rangle$, where $\{\dots\}_{\rm sym}$ is the average over the six orderings,
\begin{align}\label{eq:A2d}
S^{abc}(\v x,\v y,\v z)=d^{abc}\,\frac{N_c}{(N_c^2-4)(N_c^2-1)}\,o(\v x,\v y,\v z)\,,\qquad o(\v x,\v y,\v z)\equiv\big\langle d^{efg}\{\rho^e(\v x)\rho^f(\v y)\rho^g(\v z)\}_{\rm sym}\big\rangle\,.
\end{align}
\end{itemize}}}\par\smallskip
\noindent and the paragraph ``(A2) holds for a Gaussian projectile \dots evolution terms.'' by:\par\smallskip
\noindent\fbox{\parbox{0.97\textwidth}{\small
The normalisation follows from $d^{abc}d^{abc}=(N_c^2-4)(N_c^2-1)/N_c$; $o$ is real and totally symmetric in its arguments. By (A1), $S^{abc}=d^{abc}g_d+f^{abc}g_f$ with $g_d$ totally symmetric and $g_f$ totally antisymmetric in the three points; (A2) states that $g_f=0$ and gives $g_d$ the normalisation of Eq.~\eqref{eq:A2d}. The function $o$ is C-odd: it is the three-point function of the cubic Casimir density, i.e.\ the odderon of the projectile. It vanishes for a Gaussian projectile (Gaussian Weyl symbol of the density matrix, as in the McLerran--Venugopalan model), whose odd moments vanish, and it is non-zero, for instance, for the weight functional with a cubic term $d_{abc}\rho^a\rho^b\rho^c$ of~\cite{Jeon:2005cf}, where it is local. For $SU(2)$, $d^{abc}=0$ and (A2) reduces to $S^{abc}=0$. Nothing else is assumed; in particular no assumption on the four-point function is needed, since the four-$\rho$ rows contribute only to the evolution terms.
}}\par\smallskip
\noindent [New bib entry needed: \texttt{Jeon:2005cf} = S.~Jeon, R.~Venugopalan, \emph{A classical Odderon in QCD at high energies}, Phys.\ Rev.\ D71 (2005) 125003, hep-ph/0503219. The key is our choice; please check against INSPIRE.]

\subsubsection*{(P2) Paper, paragraph ``Quadratic and cubic terms'': the sentences after Eq.~(cubicavg)}
\noindent Replace ``Under (A2) the first term vanishes, \emph{but the contact terms do not}: \dots\ with the point $q$ replaced by $p$.'' by:\par\smallskip
\noindent\fbox{\parbox{0.97\textwidth}{\small
By (A2) the first term is $C^{abc}S^{abc}=\frac{N_c}{(N_c^2-4)(N_c^2-1)}\,C^{abc}d^{abc}\,o(p_1,p_2,p_3)$, and $2C^{abc}S^{abc}$ in the two forms with an anticommutator; we write it as $\frac{N_c}{(N_c^2-4)(N_c^2-1)}\Delta\,o(p_1,p_2,p_3)$, where $\Delta\equiv\sum C^{abc}_{\rm seq}d^{abc}$ is summed over the orderings into which the anticommutators expand. It vanishes for a Gaussian projectile, \emph{but the contact terms do not}: they are two-point terms $\propto\delta^{(2)}(p_i-p_j)\,n$, of the same order as the quadratic terms, and they are not removed by colour neutrality. \emph{The cubic terms therefore do not vanish after the average over a Gaussian projectile; only their symmetrised part does.} The contact terms are written below as quadratic terms at coincident points, $\mathfrak K|_{q\to p}$ denoting the kernel with the point $q$ replaced by $p$; the symmetrised parts are collected in Eq.~\eqref{eq:sodd}.
}}

\subsubsection*{(P3) Paper, paragraph ``Parity in $\v k$'': the last three sentences}
\noindent Replace ``The symmetrised part of the cubic terms is odd under $\v k\to-\v k$, \dots\ but not for the commutator remainders of the ordered products.'' by:\par\smallskip
\noindent\fbox{\parbox{0.97\textwidth}{\small
The symmetrised part of the cubic terms is odd under $\v k\to-\v k$, their contact part is even, and the quadratic terms are even. The part of the finite terms that is odd in $\v k$ is therefore exactly the symmetrised part, Eq.~\eqref{eq:sodd}; under (A2) it is proportional to the projectile odderon $o$, and it vanishes for a Gaussian projectile. This refines the statement that the cubic, ``first saturation'' terms are odd in $\v k$~\cite{McLerran:2016snu,Kovchegov:2018jun}: it holds for their symmetrised part, configuration by configuration, but not for the commutator remainders of the ordered products; and the symmetrised part survives the average over a projectile with $o\neq0$.
}}

\subsubsection*{(P4) Paper, paragraph ``Colour-singlet structures'': append after the coincident-point identities of $\mf D,\mf F,\mf Q$}
\noindent\fbox{\parbox{0.97\textwidth}{\small
The symmetrised parts of the cubic terms need, with $(D^e)_{bc}\equiv d^{ebc}$,
\begin{align}\label{eq:GQd}
\mf G(a,b,c)&\equiv f^{ijk}d^{lmn}\,U^{il}(a)\,U^{jm}(b)\,U^{kn}(c)\,,\qquad
\mf Q_d(a,b;c,e)\equiv i\,{\rm Tr}\{T^hU^\dagger(a)U(b)D^hU^\dagger(c)U(e)\}\,.
\end{align}
$\mf G$ is real and totally antisymmetric, $\mf Q_d$ is real, $\mf Q_d(a,b;c,e)=-\mf Q_d(e,c;b,a)$, $\mf Q_d(a,a;c,e)=\mf Q_d(a,b;c,c)=0$ and $\mf Q_d(a,b;c,a)=-\mf G(a,b,c)$. In terms of the fundamental Wilson lines $V$, $U^{ab}=2{\rm tr}(t^aVt^bV^\dagger)$, a Fierz rearrangement gives, exactly,
\begin{align}\label{eq:Gfund}
\mf G(a,b,c)&={\rm Im}\big[{\rm tr}(V^\dagger_aV_b)\,{\rm tr}(V^\dagger_bV_c)\,{\rm tr}(V^\dagger_cV_a)\big]+{\rm Im}\,{\rm tr}\big(V^\dagger_aV_bV^\dagger_cV_aV^\dagger_bV_c\big)\,,\notag\\
\mf Q_d(a,b;c,e)&=-{\rm Im}\big[{\rm tr}A\,{\rm tr}B\,{\rm tr}(A^\dagger B^\dagger)\big]-{\rm Im}\,{\rm tr}\big(ABA^\dagger B^\dagger\big)\,,\qquad A=V^\dagger_aV_b\,,\ B=V^\dagger_cV_e\,.
\end{align}
These are C-odd (odderon) structures of the target: under $V\to V^*$ they change sign, while $\mf D$, $\mf F$, $\mf Q$ do not; for a weak target they start at third order in the target field, and they vanish for $SU(2)$.
}}

\subsubsection*{(P5) Paper: new paragraph after Eq.~(sIIIVI) and the sentence ``Rows~III.V and~III.VI are cubic \dots'' (before Table~\ref{tab:singletrows})}
\noindent\fbox{\parbox{0.97\textwidth}{\small
\paragraph{Symmetrised part of the cubic rows.} The colour factors $\Delta$ of the symmetrised parts are combinations of $\mf G$ and, in Rows~II.I and~III.II, of $\mf Q_d$. With the row prefactors and kernels of the contact forms above, the part of the finite terms odd in $\v k$ is
\begin{align}\label{eq:sodd}
\frac{d^3N_{\rm odd}}{d^3k}\bigg|_{\rm singlet}&=\frac{\mathcal C_{\rm f}}{N_c^2-1}\,\frac{N_c}{N_c^2-4}\Bigg\{\int_{w,w'}e^{-ik\cdot(w-w')}\int_{x,y,x'}\frac{1}{2\pi}\,\mathfrak K^{\rm II}_{2b}\Big[\mf G(x,y,x')-\mf G(x,y,w')\Big]\,o(x,y,x')\notag\\
&\quad+\int_{w,w'}e^{-ik\cdot(w-w')}\int_{x,y,z,x'}2i\Big\{\mathfrak K^{\rm III}_z\Big[\mf G(x,y,x')-\mf G(x,y,w')+\mf G(x,x',z)+\mf G(x,z,w')\Big]\notag\\
&\qquad\qquad+\mathfrak K^{\rm III}_w\Big[\mf G(x,y,x')-\mf G(x,y,w')+\mf G(x,x',w)+\mf G(x,w,w')\Big]\Big\}\,o(x,y,x')\notag\\
&\quad-\int_{x,z,y,y',x',w'}e^{-ik\cdot(z-w')}\,i\,\mathfrak K^{\rm V}_1\Big[\mf G(x,x',z)+\mf G(x,z,w')+\mf G(y,x',y')+\mf G(y,y',w')\Big]\,o(y,x',y')\notag\\
&\quad-\int_{x,y,x',z,w,w'}e^{-ik\cdot(w-w')}\,2i\,\mathfrak K_{\mathcal T}(w,w')\Big[\mf G(x,y,x')+\mf G(x,x',z)+\mf Q_d(z,y;x',w)\Big]\,o(x,y,x')\notag\\
&\quad-\int_{w,x',y',z,y,x}e^{-ik\cdot(w-y')}\,2i\,\Big\{\Big[-\mathfrak K_{\widehat{\mathcal E}}+\mathfrak K_\Psi+\mathfrak K^{(y')}_\delta\Big]\Big[\mf G(x,y,x')-\mf G(x,y,y')-\mf G(x',z,w)+\mf G(y',z,w)\Big]\notag\\
&\qquad\qquad+\mathfrak K^{(z)}_\delta\Big[\mf G(x,y,x')-\mf G(x,y,z)-\mf G(x',z,w)\Big]\Big\}\,o(x,y,x')\notag\\
&\quad-\int_{w,y',z,y,x,x'}e^{-ik\cdot(w-y')}\,2i\,\Big\{\Big[-\mathfrak K_{\widehat{\mathcal E}}+\mathfrak K_\Psi+\mathfrak K^{(y')}_\delta\Big]\Big[-\mf G(x,y,x')+\mf G(x,y,y')-\mf G(x,x',w)+\mf G(x,y',w)\Big]\notag\\
&\qquad\qquad+\mathfrak K^{(z)}_\delta\Big[-\mf G(x,y,x')+\mf G(x,y,z)-\mf G(x,x',w)+\mf G(x,z,w)\Big]\Big\}\,o(x,y,x')\notag\\
&\quad-\int_{w,y',z,y,x,x'}e^{-ik\cdot(w-y')}\,2i\,\Big\{\Big[-\mathfrak K_{\widehat{\mathcal E}}+\mathfrak K_\Psi+\mathfrak K^{(y')}_\delta\Big]\Big[-\mf G(x,y,x')+\mf G(x,y,y')+\mf G(y,x',z)-\mf G(y,y',z)\Big]\notag\\
&\qquad\qquad+\mathfrak K^{(z)}_\delta\Big[-\mf G(x,y,x')+\mf G(x,y,z)+\mf G(y,x',z)\Big]\Big\}\,o(x,y,x')\Bigg\}+{\rm c.c.}\notag\\
&\quad-\frac{\mathcal C_{\rm f}}{N_c^2-1}\,\frac{N_c}{N_c^2-4}\int_{x,y,z,w,x',w'}4\sin\big(\v k\cdot(\v w-\v w')\big)\,\mathfrak K_{\mathcal T}(w,w')\notag\\
&\qquad\qquad\times\Big[\mf G(x,z,w')-\mf G(x,y,x')-\mf Q_d(z,y;x',w)\Big]\,o(x,y,x')\,.
\end{align}
The lines are, in this order, Rows~I.II ($N_{2b}$), I.III, I.V (first part), III.II, III.III (first part), III.V and III.VI, each with its complex conjugate, and the first part of Row~II.I ($N_{1a}+N_{1b}$, after $\v w\leftrightarrow\v w'$ in $N_{1b}$). The pairs $\mathfrak K^{(y')}_\delta$, $\mathfrak K^{(z)}_\delta$ stay under one $\v y'$ integral. None of the thirteen colour factors vanishes identically. Each term is real and odd under $\v k\to-\v k$; it is C-odd in the projectile ($o$) and in the target ($\mf G$, $\mf Q_d$), so that it vanishes after an average over a C-even target ensemble and survives configuration by configuration or for a target with an odderon component; since all kernels are invariant under the simultaneous rotation by $\pi$ of all coordinates and of $\v k$, it also integrates to zero if the averaged product of the colour factor and $o$ is invariant under the reflection of all points through a common centre. The soft-remainder pairs of Section~\ref{sec:fin_soft} pair in Eq.~\eqref{eq:sodd} as in the contact forms, since Eqs.~\eqref{eq:colid} are operator identities. The colour factors were obtained by contracting the colour structures of Appendix~\ref{app:finite} with $d^{abc}$ and verified for $SU(3)$, $SU(4)$, $SU(5)$ (to $2\times10^{-15}$), and the complete average (contact terms and symmetrised part) at the operator level with $SU(3)$ charges in a C-odd, colour-invariant state, for which (A2) holds with $o\neq0$ (to $8\times10^{-16}$).
}}

\subsubsection*{(P6) Paper, Table~\ref{tab:singletrows}: column header and caption}
\noindent Header ``after (A1), (A2)'' $\to$ ``after (A1), (A2): contact terms''; append a column ``odd part'' with entries: I.II $N_{2b}$, I.III, I.V first part, II.I first part, III.II, III.III first part, III.V, III.VI: \eqref{eq:sodd}; quadratic rows: ---. Caption:\par\smallskip
\noindent\fbox{\parbox{0.97\textwidth}{\small
Content of the rows after the average over the projectile. For every cubic row the contact terms are even in $\v k$, real after the complex conjugate is added, and do not vanish for a Gaussian, colour-neutral projectile; the symmetrised part is odd in $\v k$, proportional to the projectile odderon $o$ of Eq.~\eqref{eq:A2d}, and given in Eq.~\eqref{eq:sodd}.
}}

\subsubsection*{(P7) Paper, Remarks (iii) after the table and the sentence in Sec.~\ref{sec:fin_summary}}
\noindent Remark (iii) $\to$\par\smallskip
\noindent\fbox{\parbox{0.97\textwidth}{\small
(iii) A chiral part $f^{abc}g_f$ of the symmetrised three-point function, excluded by (A2), would add to each cubic structure $t\,\Phi\,g_f(p_1,p_2,p_3)$, with $\Phi$ the factor of the contact terms evaluated at separate points and $t=1$ ($2$) for the ordered (anticommutator) forms; it is also odd in $\v k$. For a local $o$, as in~\cite{Jeon:2005cf}, the integrability of the kernels at the triple coincidence of the charge positions has not been examined.
}}\par\smallskip
\noindent In Sec.~\ref{sec:fin_summary} replace ``After the average over a Gaussian, colour-singlet projectile the finite terms are even in $\v k$; the three-$\rho$ rows contribute through the contact terms of the commutators of the ordered charges.'' by\par\smallskip
\noindent\fbox{\parbox{0.97\textwidth}{\small
After the average over a colour-singlet projectile the three-$\rho$ rows contribute through the contact terms of the commutators of the ordered charges, which are even in $\v k$, and through their symmetrised part, Eq.~\eqref{eq:sodd}, which is odd in $\v k$ and proportional to the projectile odderon $o$; for a Gaussian projectile $o=0$ and the finite terms are even in $\v k$.
}}

\subsubsection*{(P8) Paper, \texttt{sec\_checks.tex}: paragraphs ``Counting of charges'' and ``Odd terms''}
\noindent ``Counting of charges'', last two sentences $\to$\par\smallskip
\noindent\fbox{\parbox{0.97\textwidth}{\small
The finite terms are of order $g^4\rho^2$ and $g^4\rho^3$. After the average over the projectile the $g^4\rho^3$ terms do not drop out: their symmetrised part is proportional to the symmetrised three-point function, i.e.\ to the projectile odderon $o$ of Eq.~\eqref{eq:A2d}, and the commutators produced by the reordering of the charges leave contact terms of order $g^4\langle\rho\rho\rangle$, of the same order as the quadratic finite terms (Section~\ref{sec:fin_singlet}). The three-$\rho$ terms that multiply a rapidity logarithm cancel in the symmetrised sense for any three-point function, so that no logarithm multiplies $o$.
}}\par\smallskip
\noindent ``Odd terms'' $\to$\par\smallskip
\noindent\fbox{\parbox{0.97\textwidth}{\small
\paragraph{Odd terms.} The three-$\rho$ terms carry an explicit factor $i$ (in Row~I.II, $N_{2b}$, it is in the Fourier transform of an odd kernel), and all kernels depend on $\v k$ through functions with $f(-s)=f(s)^*$. A three-$\rho$ term plus its complex conjugate is therefore odd under $\v k\to-\v k$ in its totally symmetrised part and even in the part produced by the commutators of the ordered charges (Section~\ref{sec:fin_singlet}). For a globally colour-invariant projectile with the symmetrised three-point function~\eqref{eq:A2d} the odd part is Eq.~\eqref{eq:sodd}: it is proportional to the projectile odderon $o$ and contains the target only through the C-odd structures $\mf G$ and $\mf Q_d$, Eq.~\eqref{eq:Gfund}, the imaginary parts of products of three dipoles and of sextupoles. It therefore requires a C-odd three-point function of the projectile, it vanishes after an average over a C-even target ensemble, and for a weak target it starts with three gluon exchanges; for a Gaussian (McLerran--Venugopalan-type) projectile, $o=0$, the spectrum is even in $\v k$. This agrees with the classical-field result that the leading-order spectrum is even configuration by configuration and that the odd part comes from the term cubic in the projectile charge, averages out for a Gaussian projectile and is ``defined by odderon exchanges''~\cite{McLerran:2016snu}, and with the counting of~\cite{Kovchegov:2018jun}, three scatterings in the projectile and three in the target. The even commutator part does \emph{not} vanish: the three-$\rho$ rows contribute to the averaged spectrum through two-point contact terms $\propto\delta^{(2)}(\v x_i-\v x_j)\,n$, which have no counterpart for classical charges. All colour factors of the averaged terms, the symmetrised parts and the parity statements were verified at the operator level, with charges represented by generators on two or three sites and globally colour-invariant density matrices, including a C-odd state with $o\neq0$ (Section~\ref{sec:fin_singlet}). A comparison of the three-$\rho$ terms with the dilute limit of the first saturation correction to single inclusive production of Li and Skokov~\cite{Li:2021ntm}, and a term-by-term comparison of Eq.~\eqref{eq:sodd} with the odd part of~\cite{McLerran:2016snu}, have not been carried out.
}}

\subsubsection*{(P9) Paper, \texttt{sec\_rapidity.tex}, Sec.~\ref{sec:threerho}: append to the paragraph ``For each of the two logarithms \dots''}
\noindent\fbox{\parbox{0.97\textwidth}{\small
Since the colour tensor itself vanishes, the cancellation holds for any symmetrised three-point function of the projectile, in particular for the $d^{abc}$ (odderon) one of Eq.~\eqref{eq:A2d}, which we checked separately with $SU(3)$ and $SU(4)$ matrices: no rapidity logarithm multiplies the three-point function. This is required, since the LO spectrum depends on the projectile only through the C-even two-point function $n$, and its one-step evolution cannot produce the C-odd $o$.
}}

\subsubsection*{(P10) Paper, conclusions and abstract}
\noindent \texttt{sec\_conclusions.tex}, replace ``In the averaged form the terms cubic in the charges do not vanish for a Gaussian projectile: only their symmetrised part, which is odd in $\v k$, does, and the commutators of the ordered charges leave contact terms of the same order as the quadratic terms. For such a projectile the spectrum is even in $\v k$.'' by\par\smallskip
\noindent\fbox{\parbox{0.97\textwidth}{\small
In the averaged form the terms cubic in the charges contribute in two ways. Their symmetrised part is proportional to the symmetrised three-point function of the projectile, which for a colour-singlet projectile is the $d^{abc}$ (odderon) correlator $o$; it is odd in $\v k$, and its colour factors are the C-odd target structures ${\rm Im}[{\rm tr}\,{\rm tr}\,{\rm tr}]$ and ${\rm Im}\,{\rm tr}$ of three dipoles and of a sextupole. The commutators of the ordered charges leave contact terms of the same order as the quadratic terms, which are even in $\v k$ and do not vanish for a Gaussian projectile. For a Gaussian projectile, or after an average over a C-even target, the spectrum is even in $\v k$. No rapidity logarithm multiplies the three-point function.
}}\par\smallskip
\noindent Abstract, replace the last sentence by\par\smallskip
\noindent\fbox{\parbox{0.97\textwidth}{\small
The terms cubic in $\rho$ do not vanish after the average over the projectile: their symmetrised part, odd in the transverse momentum, is proportional to the projectile odderon and to C-odd target structures, and vanishes only for a Gaussian projectile, while the commutators of the charges leave contact terms of the same order as the quadratic terms.
}}\par\smallskip
\noindent Introduction, last bullet, replace ``After the average over a Gaussian projectile the terms cubic in $\rho$ do not vanish: only their symmetrised part, which is odd in $\v k$, does, while the commutators of the ordered charges leave contact terms (Section~\ref{sec:finite}).'' by ``After the average over the projectile the terms cubic in $\rho$ do not vanish: their symmetrised part, odd in $\v k$, is proportional to the $d^{abc}$ (odderon) three-point function of the projectile, and the commutators of the ordered charges leave contact terms, even in $\v k$, which survive also for a Gaussian projectile (Section~\ref{sec:finite}).''\par
\noindent Optionally, \texttt{sec\_observable.tex} after Eq.~(singlet): ``Where the three-point function is needed we use Eq.~\eqref{eq:A2d}.''

\subsubsection*{(F1) \texttt{Final\_terms\_fixed}, fragment \texttt{build/src/singlet.tex}}
\noindent (a) Item (A2) of ``Projectile average: what is assumed'' $\to$ (A2$'$) of Sec.~\ref{sec:dabc-ass} of this note (with $\langle\{\dots\}_{\rm sym}\rangle$ as defined there), and the paragraph ``(A2) holds for a Gaussian projectile \dots'' $\to$ the first paragraph of Sec.~\ref{sec:dabc-ass} (``By (A1) \dots\ $SU(2)$ \dots''), ending with ``No assumption about the four-point function is needed (four-$\rho$ rows contribute only to the evolution terms).''\par
\noindent (b) ``Cubic terms'' paragraph: ``Under (A2) the first term vanishes, \emph{but the contact terms do not}:'' $\to$ ``Under (A2$'$) the first term is $\frac{N_{c}}{(N_{c}^{2}-4)(N_{c}^{2}-1)}\Delta\,o(p_{1},p_{2},p_{3})$, $\Delta\equiv\sum_{\rm orderings}C^{abc}d^{abc}$ (Sec.~``Symmetrised part'' below); it vanishes for a Gaussian projectile, \emph{but the contact terms do not}:''.\par
\noindent (c) ``Parity in $k$'', last sentence $\to$ ``The part of the finite terms odd in $k$ is exactly the $S^{abc}$ part; under (A2$'$) it is given in Sec.~``Symmetrised part'' below; it vanishes for $o=0$.''\par
\noindent (d) New subsection before ``Which rows are cubic, which quadratic'': ``\textbf{Symmetrised part of the cubic rows ($d^{abc}$)}'' = Secs.~\ref{sec:dabc-inv}, \ref{sec:dabc-table} (with the table and Eq.~\eqref{eq:dabc-odd}) and~\ref{sec:dabc-parity} of this note, verbatim (kernel names as in the fragment).\par
\noindent (e) Table ``Which rows are cubic, which quadratic'': header ``after the average (A1), (A2)'' $\to$ ``after the average (A1), (A2$'$): contact terms''; paragraph after the table: ``For every cubic row the symmetrised part ($S^{abc}$) is zero under (A2) and would be odd in $k$ otherwise;'' $\to$ ``For every cubic row the symmetrised part is $\propto o$, odd in $k$ and given in Eq.~(dabc-odd);''.\par
\noindent (f) ``Numerical checks'': add the items of Sec.~\ref{sec:dabc-num}. ``What is not settled'' (5) $\to$ ``(5) (A2$'$) is an assumption on the projectile: $g_{f}=0$ (Q1); for a local $o$ the integrability at the triple coincidence of the $\rho$ points is not checked (Q3).''
\endgroup
